\chapter{Introduction and Motivation}
\label{chap:introduction}

\section{Longitudinal Data and Patient Heterogeneity}
\label{sec:patient-heterogeneity}

The meaning of a patient's current clinical assessment depends on the history that precedes it. Two patients may receive the same clinical score at a visit, although one has remained stable for several years and the other has declined rapidly. Their similarity at that visit conceals a difference in progression that may matter for prognosis and follow-up. Longitudinal records make these distinctions observable by retaining the sequence of measurements, diagnoses, and clinical events through which the present state emerged. They provide a basis for studying both how an individual changes and why patients with the same diagnosis follow different courses~\cite{lee2021odtp}.

Patient heterogeneity consequently has a temporal dimension. Observed differences may reflect distinct patterns of disease, different positions along a shared course, or variation in the rate of progression. These possibilities imply different meanings for a patient subgroup. SuStaIn makes the distinction between subtype and stage explicit by jointly estimating progression patterns and patients' positions within them~\cite{young2018sustain}. Although its original formulation learns from cross-sectional imaging, the distinction is equally important when interpreting longitudinal groups: patients may look different because they are progressing differently or because they are observed at different stages. A useful temporal representation should make these alternatives easier to investigate.

\section{Learning from Irregular Clinical Histories}

The opportunity provided by longitudinal data comes with an observation problem. Clinical records contain samples of an evolving process, taken at times determined partly by care delivery. Visits are unevenly spaced, different variables are measured on different schedules, and some patients have much longer recorded histories than others. A value recorded yesterday and the same value recorded months earlier need not carry the same information about the present. Moreover, the decision to order a test may itself reflect concern about deterioration. As Hripcsak and Albers emphasize, the EHR reflects the recording process as well as the patient's condition~\cite{hripcsak2013phenotyping}.

Temporal learning therefore requires a representation of what was measured, when it was measured, and which information remains unavailable. Time-aware LSTM networks incorporate elapsed time into recurrent memory, while GRU-D uses observation masks and time intervals to learn how missing values and past hidden states should influence prediction~\cite{baytas2017tlstm,che2018grud}. These designs illustrate why temporal order alone is an incomplete description of a clinical history. They also expose an interpretive challenge: similarities in a learned representation can arise from the course of disease, the pattern of observation, or both. Handling irregularity therefore shapes both the information available for prediction and the comparisons that can be made between patients.

\section{From Temporal Representations to Clinical Prediction}

A representation of patient history becomes useful for a particular prediction problem through its learning objective. The objective may be regression, such as estimating a future clinical score or remaining length of stay, or classification, such as predicting deterioration or diagnostic status. Harutyunyan et al.'s clinical time-series benchmark illustrates this range through length-of-stay, mortality, decompensation, and phenotype-classification tasks~\cite{harutyunyan2019benchmark}. When the target concerns the timing of an event and follow-up may end before that event occurs, survival analysis uses both observed events and the partial information supplied by censoring. The target and observation scheme determine what constitutes a valid prediction and which information the representation is encouraged to retain.

These predictions can be revised as the patient's history grows. Statistical joint models established a framework for updating survival probabilities from repeated marker measurements~\cite{rizopoulos2011dynamic}; Dynamic-DeepHit develops a neural approach that learns event-time distributions from longitudinal histories while accommodating competing risks~\cite{lee2020dynamicdeephit}. The shared clinical objective is to keep prognosis responsive to the evidence available at each assessment. In this sense, dynamic prediction concerns the evolving information about a patient, without requiring the model parameters to be retrained after every observation.

Individual prediction, however, leaves a further question about the population. An accurate estimate of an endpoint does not describe whether patients reach that endpoint through similar courses. Likewise, two histories assigned the same predicted risk may differ in clinically important ways. This motivates learning a population structure alongside prognosis: a set of groups through which individual histories can be compared, summarized, and interpreted.

\section{From Temporal Clustering to Supervised Phenotype Discovery}

Temporal clustering supplies such a structure by grouping sequences according to a chosen notion of similarity. A natural starting point is similarity in the observations themselves, expressed through trajectory distances or a learned representation. In the patient-subtyping approach of Baytas et al., a time-aware autoencoder represents irregular records before clustering patients into subtypes~\cite{baytas2017tlstm}. This makes recurring temporal structure accessible without specifying an outcome in advance. Yet the variation required to reconstruct a history need not coincide with the variation most relevant to a clinical question. A clustering objective must therefore express what the groups are intended to preserve.

Supervised temporal clustering addresses this choice by allowing clinical outcomes to guide group formation. The supervision can come from the same target families used for individual prediction. AC-TPC learns discrete representations that approximate the outcome distributions associated with patient histories, and its supplementary formulation also accommodates continuous regression targets~\cite{lee2020actpc}. OG-DTC provides a clinical example in which temporal reconstruction and clustering are jointly optimized with outcome regression~\cite{wang2024ogdtc}. For categorical outcomes, CAMELOT learns cluster representations through which predictions are made, with additional objectives addressing class imbalance and cluster collapse~\cite{aguiar2022camelot}. Here, outcome supervision gives the groups an explicit clinical purpose while their assignments remain quantities to be learned.

Event-time outcomes provide a further source of supervision. ODTP learns temporal phenotypes that summarize the expected type and timing of subsequent events~\cite{lee2021odtp}. VaDeSC-EHR instead integrates diagnosis-sequence reconstruction, latent clustering, and survival information within a generative model~\cite{qiu2025vadescehr}. These formulations give different answers to the same design question: whether patient similarity should primarily preserve a predictive distribution or the joint structure of clinical histories and outcomes. The distinction matters because a group can be prognostically useful while retaining substantial heterogeneity in the histories of its members.

T-Phenotype responds directly to this tension by associating outcome similarity with predictive temporal patterns, using path-based connectivity to distinguish groups that a comparison of predicted outcomes alone may merge~\cite{qin2023tphenotype}. These works make supervised temporal clustering a substantive choice about which aspects of heterogeneity to retain. A representation optimized for reconstruction, outcome prediction, or both can lead to different groups. Each grouping must therefore be assessed through the aspects of progression it captures and the evidence supporting their clinical interpretation.

In clinical research, interpreting these groups as \emph{phenotypes} gives the grouping a substantive meaning. EHR phenotyping broadly includes deriving clinically relevant characteristics and identifying predefined conditions, as well as discovering patient subgroups~\cite{hripcsak2013phenotyping}. This report concerns temporal phenotype discovery: identifying groups characterized by recurring longitudinal patterns and their clinical associations. Outcome labels help define those associations, but they do not supply ground-truth phenotype assignments. A phenotype may distinguish patients within an outcome class, and its interpretation must be supported by the characteristics and reproducibility of the resulting group.

\section{Why Temporal Phenotyping Needs to Be Dynamic}

A group discovered from complete trajectories offers a retrospective description of a disease course. Clinical follow-up presents a different setting: the patient's future trajectory is still unknown, and the information available for assignment changes from visit to visit. Returning to the patients with similar clinical scores, an initially stable history may later show sustained deterioration. A useful assessment should be able to reflect that new evidence. Assigning a phenotype to the history available at each assessment allows the patient to be compared with previously observed histories that are relevant to the current course.

AC-TPC and ODTP implement this perspective through phenotype assignments that can change as observations accumulate~\cite{lee2020actpc,lee2021odtp}. The clinical motivation is to describe an evolving prognosis in relation to recognizable population patterns. Such descriptions could help characterize when a patient's course begins to resemble a different prognostic group and which recorded changes accompany that shift. Their value lies in combining the comparative information supplied by a cohort with an assessment responsive to the individual patient's history.

The interpretation of change depends on the phenotype being modeled. A progression subtype may describe a relatively persistent course, whereas a temporal phenotype may describe a changing clinical or prognostic state within that course. The subtype--stage distinction highlighted by SuStaIn helps clarify why a changing assignment need not imply a change in disease subtype~\cite{young2018sustain}. More generally, a transition can reflect new clinical evidence, revised uncertainty, or a change in observation patterns. Dynamic phenotyping must therefore balance responsiveness with reproducibility: frequent reassignment can obscure meaningful progression, while a fixed assignment can conceal important changes. Assessing that balance requires examining the trajectories and outcomes associated with each group, the uncertainty of assignment, and the information available when the assignment was made.

\section{Motivation and Scope of the Review}

The resulting problem is to learn a clinically informative description of patient heterogeneity from incomplete histories while preserving the ability to update that description over time. This report asks: \emph{What should a temporal phenotype preserve about a patient's progression and prognosis, and how can its assignment remain meaningful as new observations become available?} Addressing this question requires examining the relationship between the representation, the learning objective, and the interpretation supported by the resulting groups.

The four assigned works provide complementary perspectives. Dynamic-DeepHit establishes the predictive foundation by updating event-time distributions from longitudinal information. ODTP turns event-time behavior into a basis for temporal phenotype discovery. CAMELOT examines prediction through learned groups and their feature--time characterization, while VaDeSC-EHR jointly models longitudinal histories and survival. The following chapters develop a common formulation and critically examine each method's assumptions, technical construction, and empirical evidence. Their comparison will distinguish evidence for accurate prognosis from evidence that a grouping provides a reproducible and clinically informative account of disease progression.
